\chapter{Detailed Architecture Diagrams}
\label{append:architecture}
\vspace*{2em}

This appendix reproduces the system's master architecture diagram, from which
the simplified figures in \autoref{ch:design} and \autoref{ch:implementation}
were drawn. It shows one system at three magnifications: \emph{View~1}
(deployment) shows where the system runs; one box in it, the backend service,
is the whole of \emph{View~2} (application components), which shows how the
modules connect; and one module in that, the chat handler, is the whole of
\emph{View~3} (the AI agent), which traces a single chat turn.

\autoref{append:fig:arch-all} overlays four use cases on the diagram, each in
its own colour; \autoref{append:fig:arch-s1}--\autoref{append:fig:arch-i2}
show them one at a time, with the components on each path ringed and the path
numbered in order:

\begin{itemize}[nosep]
  \item \textbf{S1} (orange): a student asks the assistant about course material;
  \item \textbf{S2} (blue): a student takes an adaptive quiz;
  \item \textbf{I1} (red): an instructor requests a course analysis;
  \item \textbf{I2} (teal): an instructor connects Canvas and pulls courses.
\end{itemize}

Two differences from the system as evaluated should be kept in mind. First,
View~1 shows the \emph{production target} on AWS (managed database, container
orchestration, content delivery network, and managed identity), which has
been designed but not provisioned; the deployed prototype runs on a single
EC2 host (\autoref{design:fig:architecture}). Second, the diagram predates the
move to hosted inference on Groq: it shows Gemini as the managed model
provider, which has since been replaced by Groq
(\autoref{impl:subsec:llm-providers}).

\begin{figure}[p]
  \centering
  \includegraphics[height=0.86\textheight,keepaspectratio]{\RootPicDir/appendix/arch-all-use-cases.pdf}
  \caption{StudyAI master architecture: deployment (View~1), application
           components (View~2), and the AI agent (View~3), with all four use
           cases highlighted.}
  \label{append:fig:arch-all}
\end{figure}

\begin{figure}[p]
  \centering
  \includegraphics[height=0.86\textheight,keepaspectratio]{\RootPicDir/appendix/arch-s1-ask-ta.pdf}
  \caption{Use case S1: a student asks the assistant about course material.
           The request passes through retrieval, the optional compression
           stage, the persona, and the agent loop, and the answer streams
           back token by token.}
  \label{append:fig:arch-s1}
\end{figure}

\begin{figure}[p]
  \centering
  \includegraphics[height=0.86\textheight,keepaspectratio]{\RootPicDir/appendix/arch-s2-adaptive-quiz.pdf}
  \caption{Use case S2: a student takes an adaptive quiz. The fuzzy engine
           sets the target Bloom's level and complexity; a question is taken
           from the bank or, on a miss, generated by the model and cached.}
  \label{append:fig:arch-s2}
\end{figure}

\begin{figure}[p]
  \centering
  \includegraphics[height=0.86\textheight,keepaspectratio]{\RootPicDir/appendix/arch-i1-analyse-course.pdf}
  \caption{Use case I1: an instructor requests a course analysis. Class
           signals are recomputed from the database and combined with live
           web-search results in one prompt; the note streams back.}
  \label{append:fig:arch-i1}
\end{figure}

\begin{figure}[p]
  \centering
  \includegraphics[height=0.86\textheight,keepaspectratio]{\RootPicDir/appendix/arch-i2-connect-lms.pdf}
  \caption{Use case I2: an instructor connects Canvas and pulls courses. The
           OAuth token is stored per user, and courses, assignments, and
           grades are fetched and normalised; this path involves no AI.}
  \label{append:fig:arch-i2}
\end{figure}
